В работе исследован переход от локальной длительности задачи с ИИ к сроку
фиксированного набора связанных программных работ. Для одной
человеко-агентной ячейки ответ на поставленный во введении вопрос состоит в
следующем: конфигурация сокращает срок относительно последовательной работы без
ИИ не в силу самого наличия нескольких агентов, а только тогда, когда
существует допустимое расписание, завершающее тот же принятый объём работ
быстрее исходного сценария. Достижимое ускорение определяется совместно
нормированными
длительностями задач, их зависимостями и фазовыми профилями, числом и правилами
использования агентных потоков, интеграционными издержками и последовательным
ресурсом человеческого внимания. Тем самым подтверждён основной тезис:
ускорение является свойством полностью заданного процесса, а не изолированной
характеристикой ИИ-инструмента или номинального числа агентов.

Первый результат работы --- операциональная рамка сравнения --- задаёт границы
такого вывода. Коэффициент $k_i^{(r)}$ описывает длительность принятого
результата для конкретной тройки ``задача --- инструмент --- режим'' и не имеет
причинной интерпретации без отдельного дизайна сравнения. Конфигурация
$\sigma=(P,\rho)$ фиксирует число потоков и назначение режимов, но корректное
сопоставление сценариев требует также сохранять либо явно изменять
декомпозицию, DAG, фазовые профили, критерий приёмки и функции издержек. Это
разделение не позволяет переносить измеренный $k$ из интерактивного режима в
делегированную конфигурацию с другим $P$ и отличает автономную агентную фазу,
во время которой человек свободен, от полностью автономного процесса.

Второй результат --- иерархия M0--M4 --- показывает, как по мере снятия
идеализаций меняется сила ответа. На M0 ускорение достигается тогда и только
тогда, когда $P>\bar{k}_w$. На M1--M2 это условие остаётся необходимым, но
неделимость и зависимости создают дополнительный зазор; при целом числе
одинаковых потоков и $C(P)=1$ достаточную гарантию жадного расписания даёт
$\bar{k}_w<P-(P-1)\mu$, где $\mu=M_N/T_h$ для M1 и $\mu=L/T_h$ для M2. На
M3 эффективные веса включают межагентные интеграционные издержки. На M4
полностью заданный сценарий ускоряет работу тогда и только тогда, когда его
оптимальный makespan удовлетворяет $T^*<T_h$. Если структурная граница
$B_4\ge T_h$, ускорение невозможно; если $B_4<T_h$, этого ещё недостаточно.
Достаточным свидетельством ускорения служит допустимое расписание
$T(\pi)<T_h$, а точный решатель устанавливает существование и оптимальность
такого расписания.

Третий результат уточняет структуру M4. Нижняя граница

\[
  B_4=\max\left\{\frac{W_4(P)}{P},\;L_4(P),\;\gamma(P)H\right\}
\]

объединяет суммарную загрузку агентных потоков, критический путь исходного
task-DAG и единственный человеческий ресурс. Однако повторные обращения задач
к разработчику и удержание ожидающими задачами агентных слотов создают
смешанные ресурсные цепочки, не содержащиеся ни в одной из трёх ветвей.
Ресурсно-дополненный фазовый граф представляет точный срок как минимальную
длину максимального пути по допустимым назначениям и ресурсным порядкам. Тем
самым строго разведены структурная граница $B_4$, доказанный оптимум $T^*$ и
срок выбранной политики $T(\pi)$. Последовательная human-work одновременно
задаёт независимый потолок
$S_E^{\mathrm{ach}}\le1/[\gamma(P)\bar h_w]\le1/\bar h_w$: после его
достижения добавление агентов не заменяет сокращение человеческих фаз,
автоматизацию проверки или более формальный критерий приёмки.

Четвёртый результат --- серия EXP-00--EXP-08 --- проверяет эти различия на
воспроизводимых сценариях. Иллюстративные варианты 2--4 получили доказанные
оптимумы, но равенство $T^*=B_4$ оказалось частным свойством: в EXP-01 строгий
разрыв найден в 30 из 90 точных сценариев. В EXP-03 разные фазовые профили при
одинаковых $A,H,W_4,L_4$ и $B_4$ дали разные $T^*$ в 14 из 24 пар. При
неизменных длительностях дополнительная доступная мощность не ухудшала оптимум,
тогда как растущие $C(P)$ или $\gamma(P)$ создавали внутренний минимум по
настроенному $P$ на исследованных кривых EXP-02. Положительные издержки
декомпозиции дали конечную оптимальную гранулярность в 12 из 13 кривых EXP-04.
В EXP-05 interaction contrast параллелизма и дробления был положительным в 96
из 100 внутренних сравнений и нулевым в четырёх, но абсолютный эффект
декомпозиции оставался отрицательным в 26 точках. Следовательно,
взаимодополняемость параллелизма и гранулярности не означает безусловной
полезности дробления.

Проверки устойчивости ограничивают силу этих выводов. Масштабная
инвариантность воспроизведена во всех точках EXP-06A, а преимущество
иллюстративного варианта 4 над вариантом 3 сохранилось в заданных случайных и
неблагоприятных семействах ошибок EXP-06B--C. Синтетический пилот EXP-07
поддержал агрегированные направления для фиксированной политики, но не
универсальные знаки для каждого DAG и не утверждения о $T^*$. Исправленный
парный дизайн EXP-08 показал, что общий объём человеческой работы задаёт
ресурсный потолок, тогда как её размещение влияет на пути только при различии
масштабов agent- и human-фаз; фазовый порядок при этом может менять очередь и
срок даже при неизменных агрегатах. Все частоты относятся к замороженным
синтетическим матрицам и не являются оценками распространённости эффектов в
реальных проектах.

Практическое применение модели поэтому начинается не с выбора максимального
числа агентов. Для сравниваемых сценариев необходимо в общей шкале оценить
задачи и режимы, зафиксировать критерий приёмки, построить DAG и фазовые
профили, вычислить $B_4$ как диагностику обязательных ограничений, а затем
построить точное или эвристическое ресурсно-допустимое расписание и проверить
чувствительность к $k$, $h$, $C$ и $\gamma$. Активная ветвь $B_4$ показывает,
что ограничивает нижнюю границу, но только достижимость границы или анализ
ресурсного пути позволяет перенести этот вывод на фактический оптимум. В этом
смысле модель предназначена прежде всего для сравнения способов организации
работы и поиска рычагов изменения срока, а не для обещания абсолютного прогноза
без локальной калибровки.

Полученный ответ ограничен принятой архитектурой: фиксированным заранее
известным DAG, одним разработчиком, однородными агентными потоками,
непрерываемыми фазами и локальным удержанием слота. Модель не оценивает качество
кода и стоимость как самостоятельные цели и не заменяет причинное исследование
эффективности ИИ. Следующие шаги --- эмпирически измерить agent/human-интервалы,
очередь, повторную работу, $C(P)$ и $\gamma(P)$ на реальных трассах; перейти к
интервальной и стохастической калибровке; а затем расширить постановку на
динамическое раскрытие DAG, адаптивную декомпозицию и другие архитектуры
оркестрации. До такой проверки главный вывод остаётся структурным, но уже
операциональным: полезный параллелизм возникает не из числа запущенных агентов,
а из согласованного проектирования задач, зависимостей, фаз, ресурсов и правил
приёмки всего человеко-агентного процесса.
